% TOP_PROBLEM: {"problemId": "problem.lions-open-problem-on-small-time-global-exact-null-controllability-of-the-navier-stokes-equations-with-dirichlet-boundary-conditions", "canonicalTitle": "Lions' Open Problem on Small-Time Global Exact Null Controllability of the Navier-Stokes Equations with Dirichlet Boundary Conditions", "releaseRank": 329, "primaryDomain": "applied_computational_mathematics", "primaryDomainLabel": "Applied & computational mathematics", "displayStatus": "open", "releaseStatus": "open"}
% !TeX program = lualatex
% ATTEMPT_STATUS: unresolved
% Model: GPT-6 Astra Ultra; continuation dated 27 September 2026.
\documentclass[11pt]{article}
\usepackage[margin=25mm]{geometry}
\usepackage{fontspec}
\setmainfont{Latin Modern Roman}
\usepackage{amsmath,amssymb,mathtools}
\usepackage{unicode-math}
\setmathfont{Latin Modern Math}
\usepackage{xurl,hyperref}
\hypersetup{colorlinks=true,urlcolor=blue}
\setcounter{secnumdepth}{0}
\setlength{\parindent}{0pt}
\setlength{\parskip}{5pt}
\setlength{\emergencystretch}{3em}
\begin{document}
\section*{\texorpdfstring{Lions' Open Problem on Small-Time Global Exact Null Controllability of the Navier-Stokes Equations with Dirichlet Boundary Conditions}{Lions' Open Problem on Small-Time Global Exact Null Controllability of the Navier-Stokes Equations with Dirichlet Boundary Conditions}}\label{329}
\subsection{Definitions and mathematical statement}
For a bounded fluid domain $\Omega$ and a nonempty relatively open boundary portion $\Gamma$, consider
\[
 \partial_tu+(u\cdot\nabla)u-\nu\Delta u+\nabla p=0,
 \qquad\nabla\cdot u=0,
\]
with viscosity $\nu>0$. The boundary control is $u=f$ on $\Gamma$ and $u=0$ on $\partial\Omega\setminus\Gamma$. The assertion seeks, for every $T>0$ and admissible divergence-free $u_0$ in the natural $L^2$ phase space, a control and weak solution with $u(0)=u_0$ and $u(T)=0$. Incompressibility requires zero total boundary flux whenever traces are defined. ``Global'' means no smallness restriction on $u_0$; ``small-time'' means every positive $T$. The catalogue leaves dimension, boundary regularity, control spaces, and weak-solution traces implicit; equations (18), (20) and Conjecture 1(2) of its source supply those analytic conventions, which cannot be replaced by an arbitrary weak-solution class.
\par\smallskip\textit{Formulation and concept references:} \href{https://link.springer.com/article/10.1007/s40324-024-00363-7}{[S1]}.

\subsection{Short English statement}
Can boundary control on any nonempty open boundary portion bring every admissible incompressible flow exactly to rest in any prescribed positive time?
\subsection{Research attempt}
\textbf{Analytic conventions and known scope.}
The boundary-control paragraph of [S1] preceding equation (18) specifies
dimensions two or three and a bounded domain with Lipschitz-continuous
boundary.  Its earlier internal-control setup uses the stronger
$W^{2,\infty}$ regularity; that is not silently imposed here as the
boundary target's hypothesis.  The source's phase spaces are
\[
 H=\{v\in L^2(\Omega)^N:\operatorname{div}v=0,
                     \ v\cdot n=0\text{ on }\partial\Omega\},\qquad
 V=H_0^1(\Omega)^N\cap H.
\]
The normal trace in the definition of $H$ is understood weakly.
Conjecture 1(2) asks for controls in $L^2(\Gamma\times(0,T))^N$ and an
associated weak state in $C^0([0,T];H)$ with terminal state exactly zero.
This attempt retains those conventions.  In particular it does not enlarge
the solution class to arbitrary distributional solutions or replace the
uncontrolled Dirichlet boundary by slip conditions.  The source records local
null controllability and large-time results, and treats control on the entire
boundary as a solved special geometry.  Those assertions do not give arbitrary
small-time global control on a prescribed proper boundary portion.

\textbf{Flux and the correct energy identity.}
For a sufficiently regular controlled solution, integration of
$\operatorname{div}u=0$ gives
\[
 \int_\Gamma f\cdot n\,dS=0.
\]
This is a necessary compatibility condition, not an energy estimate or a
sufficient condition for controllability.  In a state class with zero normal
trace at all times there are corresponding stronger trace restrictions; the
displayed identity remains necessary.  For weak boundary controls one must
use the source's weak trace framework instead of assigning pointwise boundary
values to an arbitrary $L^2$ velocity.

Multiplying the smooth equation by $u$ and integrating gives
\[
 \frac12\frac d{dt}\int_\Omega|u|^2
 +\nu\int_\Omega|\nabla u|^2
 =\int_{\partial\Omega}
 \left(\nu u\cdot\partial_nu-p\,u\cdot n
                     -\frac12|u|^2u\cdot n\right)dS.
\]
The convection term is the divergence of $|u|^2u/2$, the pressure term is
the divergence of $pu$, and integration by parts produces the viscous
boundary term.  Thus an active Dirichlet control can exchange energy with the
fluid even when its total flux is zero.  Applying the homogeneous-boundary
energy inequality to such a solution would drop precisely the boundary work
that permits control.  Conversely, the possibility of arbitrary boundary
work does not by itself show that every interior velocity can be cancelled.

\textbf{What dissipation plus local control actually yields.}
First set the control to zero.  In two dimensions, use the standard global
energy solution and its usual regularization.  If $\lambda_1>0$ is the
Poincar\'e constant, its energy inequality implies
\[
 \|u(t)\|_2^2\leq e^{-2\nu\lambda_1t}\|u_0\|_2^2,
 \qquad
 \nu\int_t^{t+1}\|\nabla u(s)\|_2^2\,ds
 \leq\frac12e^{-2\nu\lambda_1t}\|u_0\|_2^2.
\]
There is therefore a regular time $\tau\in[t,t+1]$ with
\[
 \|u(\tau)\|_{H^1}^2
 \leq\frac{1+\lambda_1^{-1}}{2\nu}
               e^{-2\nu\lambda_1t}\|u_0\|_2^2.
\]
Let $\kappa$ be the small-data threshold for the source's local null-control
theorem over a unit time interval.  Choosing $t$ large enough makes this
quantity at most $\kappa^2$.  Apply local control starting at $\tau$, reach
zero at $\tau+1$, and then maintain zero.  For every sufficiently large
prescribed terminal time $T$, one can choose the waiting window before
$T-1$ and obtain a control ending by $T$.  Concatenation uses the same state
and boundary conventions as the two constituent results.

This is a conditional derivation from standard two-dimensional energy theory
and the cited local-control theorem, not a fresh proof of that theorem.
The waiting time grows with the size of $u_0$; it cannot be removed when the
target requires an arbitrarily small $T$.  Exponential decay also never
asserts exact extinction at a finite time.  Already the scalar Stokes-mode
equation $a'+\nu\lambda a=0$ has the nonzero solution
$a(t)=a(0)e^{-\nu\lambda t}$ for every finite $t$ if $a(0)\ne0$.

\textbf{Why compressing a large-time trajectory changes the problem.}
Suppose $u,p$ solve the equation at viscosity $\nu$, and set, for $a>0$,
\[
 v(t,x)=a\,u(at,x),\qquad P(t,x)=a^2p(at,x).
\]
Both the time derivative and convection scale by $a^2$, whereas the viscous
term scales by $a$.  Direct substitution yields
\[
 \partial_tv+(v\cdot\nabla)v-\nu\Delta v+\nabla P
 =\nu a(a-1)\Delta u(at,x).
\]
The rescaled pair solves an unforced Navier--Stokes equation only if one
also changes viscosity to $a\nu$, or in a special case where the residual
vanishes.  It also changes the initial velocity from $u_0$ to $a u_0$.
Consequently a large-time control cannot be compressed into a small-time
control for the same datum and viscosity by this substitution.  A rapid
inviscid design must still control the viscous correction and the no-slip
boundary layer on the uncontrolled boundary.

\textbf{An exact finite-dimensional control calculation.}
A linear Galerkin model clarifies what controllability estimates would need
to survive passage to the PDE.  Consider
\[
 a_i'(t)+\lambda_i a_i(t)=b_i v(t),\quad 1\leq i\leq M,
 \qquad \lambda_i>0,\quad b_i\ne0,
\]
where the $\lambda_i$ are distinct and a single scalar control is used.
Let $A=\operatorname{diag}(\lambda_i)$ and $b=(b_i)$.
The terminal state is
\[
 a(T)=e^{-AT}a(0)+\int_0^T e^{-A(T-s)}b\,v(s)\,ds.
\]
The Gram matrix of the control map has entries
\[
 W_{ij}=b_ib_j\frac{1-e^{-(\lambda_i+\lambda_j)T}}
                         {\lambda_i+\lambda_j}.
\]
It is positive definite.  Indeed for a nonzero vector $z$,
$z^TWz=\int_0^T(\sum_i z_ib_i e^{-\lambda_i t})^2dt$.
If this vanished, the analytic exponential sum would vanish identically.
Its first $M$ derivatives at zero form a Vandermonde system in the distinct
$\lambda_i$, forcing $z_i b_i=0$ for every $i$, a contradiction.
Thus
\[
 v(s)=-b^Te^{-A(T-s)}W^{-1}e^{-AT}a(0)
\]
is an exact null control.  This establishes all-time control for every fixed
finite model with these specified coefficients.

Even for one mode the minimum control energy is
\[
 \inf\int_0^T|v|^2dt
 =\frac{|a(0)|^2e^{-2\lambda T}}
        {b^2(1-e^{-2\lambda T})/(2\lambda)}
 \sim\frac{|a(0)|^2}{b^2T}\qquad(T\downarrow0),
\]
by Cauchy--Schwarz, with equality for a multiple of the exponential kernel.
In larger models the smallest eigenvalue of $W$ may become very small.
Invertibility for each $M$ does not give a bound uniform in $M$, and the
true boundary coupling is not an arbitrary list of nonzero $b_i$.
The calculation therefore does not prove even the full PDE's control theorem
without an observability estimate.  For the nonlinear equation an additional
global argument is required beyond linear observability.

\textbf{The lifting reduction and its actual nonlinear terms.}
For smooth compatible boundary data, suppose $w$ is a divergence-free lifting
of that data and write $u=z+w$.  Then $z$ has homogeneous boundary values and
satisfies
\[
 \partial_tz-\nu\Delta z+(z\cdot\nabla)z
 +(w\cdot\nabla)z+(z\cdot\nabla)w+\nabla p
 =-\partial_tw+\nu\Delta w-(w\cdot\nabla)w.
\]
This identity follows by expanding each quadratic term; it is not an existence
assertion for a lifting with arbitrary $L^2$ trace.  In a local fixed-point
scheme, smallness can control the quadratic term and the mixed terms in
suitable energy spaces.  For a large rapid boundary motion, $w$ and
$\nabla w$ need not be small.  Their terms cannot be discarded merely because
the control acts at the boundary.  A proof of the target must estimate these
effects in a space supporting the prescribed weak solution and exact terminal
trace, uniformly enough to treat every initial velocity and time.

\textbf{Verification and outcome.}
The saved diagnostic computes the finite Gram matrix and the terminal control
identity numerically for a specified three-mode system, checks the one-mode
cost formula, and verifies the time-scaling residual coefficients exactly.
It does not simulate the nonlinear boundary-controlled PDE.  The results
explain large-time control from dissipation and local control, and establish
an exact finite-dimensional null-control subcase.  The unresolved step is a
global small-time construction compatible with the uncontrolled Dirichlet
boundary and the source's state class; no such construction is supplied here.
\subsection{Sources}
\begin{itemize}
\item[S1]Enrique Fernandez-Cara, Remarks on control and inverse problems for PDEs, SeMA Journal 82 (2025), Conjecture 1(2), equations (18) and (20).
\url{https://link.springer.com/article/10.1007/s40324-024-00363-7}.
\textit{Formulation source.}
\end{itemize}
\end{document}
